\typeout{NT FILE chapter-rs485-interface.tex}%
\chapter{RS485 Interface}
\label{cha:rs485-interface}

This chapter describes the \gls{rs485-interface} used in this internship.
First, the \gls{RS485} standard will be introduced, followed by a description
of the specific implementation used in this project.

\section{RS485 Standard}
\textcolor{Red}{Review}
\gls{RS485} is a standard for serial communication that defines the electrical characteristics
of the drivers and receivers used in a multi-point communication system. It is widely used in
industrial and commercial applications due to its robustness and ability to support long-distance communication.
The \gls{RS485} standard allows for up to 32 devices to be connected on a single bus, with a maximum cable length
of 1200 meters. It uses differential signaling, which helps to reduce noise and improve signal integrity over long distances
\cite{wikipediaRS485}\cite{wang2026rs485}.

\section{Objective and Motivation}
\textcolor{Red}{Review}
As mentioned above \gls{RS485} is widely used in industrial and commercial
applications due to its robustness and ability to support long-distance
communication. The objective of the interface developed in this internship is
to provide a system that can be easily integrated with a \gls{GNSS} signal
generator, allowing for remote control and monitoring of the signal generation
process. Thus, it was decided that the interface should be able to support a
flexible way to declare protocols for both transmission and reception.

\section{Design and Implementation}
The \gls{rs485-interface} was implemented through the use of an event-driven
architecture, where the system listens for incoming data on the \gls{RS485} bus
and processes it accordingly. The interface is designed so that a user may only
need to declare C++ structs that represent protocol's messages, and a
corresponding handler function that will be called when a message of that type
is received. This allows a user to easily define new protocols and handle them
without needing to modify the core of the system, thus providing a flexible way
to manage different communication needs. Unlike the \gls{sdr-control-system},
the \gls{rs485-interface} does not use \gls{gRPC} for reflection, as that would
add unnecessary overhead to the system, it istead uses the boost-pfr
library\footnote{''Boost.PFR is a C++14 library basic reflection'' It provides
    an easy to use reflection mechanism\cite{boost-pfr}} to achieve a similar
result, allowing the system to automatically determine the structure of the
messages being sent and received, without the need for manual parsing or
serialization. This design choice allows for a more efficient and streamlined
communication process, while still providing the flexibility needed to support
different protocols. In confomity with the choices made in
\gls{sdr-control-system}, the \gls{rs485-interface} also uses errors as values
and the XMake build system. Another relevant design choice was to leave the
library as basic as possible, this could be changed in the future of course,
but for now the library offers only necessary components to read, write and
register protocols, leaving more complex features such as looping reads from
the bus, or protocol composition, to be implemented by the user on top of the
library, if needed.

\section{Architecture Diagram}
\textcolor{Red}{Review}
Diagram \ref{fig:rs485-interface-architecture} illustrates the architecture of
the \gls{rs485-interface}, showing the main components and their interactions.
This type of architecture allows for a clear separation of concerns, with each
component responsible for a specific aspect of the system, and provides clear
way for a register new protocols and handlers without needing to modify the
core of the system, thus providing a flexible way to manage different
communication needs.

\begin{figure}[htbp]
    \centering
    \input{2-MainMatter/rs485-interface-architecture.tex}
    \caption{RS485 Interface Architecture}
    \label{fig:rs485-interface-architecture}
\end{figure}

\section{Functionality and Use Cases}
\textcolor{Red}{TODO}